\scanpage{18}
\mspage{32}

\uline{Corollaire} \deleted{3}\handwritten{4}.13. Soit $A\handwritten{\to}B$ un homomorphisme local d'anneaux locaux noethériens, on suppose $k(B)/k(A)$ à multiplicité radicielle finie (par exemple une extension séparable, ou de type fini)\handwritten{,} \deleted{Conditions}\handwritten{\inserted{\uline{ou} $\underline{r}(B)=\underline{r}(A)B$.}} On munit $A$ et $B$ de leurs topologies habituelles. Alors les conditions suivantes sont \deleted{\unsure{e|é}}équivalentes: (i) $B$ est formellement simple sur $A$ (ii) $B$ est plat sur $A$, et $B_{\unsure{o|0}}=B\otimes_A k\deleted{\unsure{A|(}}$ est une algèbre universellement régulière sur $k$ (cf 2.6.) .

En effet, les conditions (i) et (ii) étant invariantes par changement de $B$ en $\overset{\handwritten{\deleted{\unsure{\wedge|\sim}}\,\wedge}}{B}$, on est ramené au cas où $B$ est complet, où cela est contenu dans \deleted{3}\handwritten{4}.11. \handwritten{\inserted{resp. 3.12.}} compte tenu que pour $B_o$, formellement simple = \unsure{universellemen|universellement} régulier en vertu de \deleted{2}\handwritten{3}.6.

Du point de vue pratique, les résultats du présent numéro consistent essentiellement en \deleted{3}\handwritten{4}.13. \deleted{l'existence}, \deleted{et} l'\uline{unicité} \deleted{de $B$} (moyennant les conditions de \deleted{3}\handwritten{4}.13.) du $B$ correspondant à un $B_o$ donné, sous réserve que $B$ soit complet pour la topologie $\underline{r}(A)$-adique (cette \deleted{ceci} unicité étant un corollaire immédiat d'ailleurs, puisqu'on dispose \deleted{\unsure{p|b}}pour $B$ des deux propriétés complémentaires formellement simple + plat), \inserted{et enfin l'\uline{existence} moyennant la séparabilité de $K/k$ .} A force de dévissage et de généralité, cette substance n'apparait pas assez clairement, et la terminologie \textquotedbl{}algèbre de Cohen\textquotedbl{} introduite concur\deleted{\unsure{r|e}}emment à \textquotedbl{}formellement simple\textquotedbl{} risque \deleted{peut-être} d'embrouiller \handwritten{\inserted{plutôt}} les idées (d'autant plus que nous n'utiliserons plus jamais que \textquotedbl{}formellement simple\textquotedbl{} ). Aussi le présent numéro rendra peut-être plus service si on commence à énoncer simplement \deleted{3}\handwritten{4}.13. en théorème, avec le corollaire bien senti précédent, et qu'on
% slanted
\margin{\handwritten{(garder cependant 4.1. \uuline{au N\textsuperscript{o}1})}}lesdémontre simultanément en deux pages ou trois, \deleted{avec} en laissant de côté tout le fatras anneaux admissibles, algèbres de Cohen etc.
% slanted
\margin{\handwritten{On peut signaler 4.6. en remarque, pour réserver l'avenir.}}Commencer par l'existence d'un $B$ satisfaisant à la fois (i) et (ii) si $K/k$ est séparable, se ramener à ce cas pour \deleted{3}\handwritten{4}.13. par descente comme dans \deleted{3}\handwritten{4}.11. prouver \deleted{3}\handwritten{4}.13. alors par comparaison du $B$-type avec le $B$ donné, ce qui donne en même temps l'unicité, et un point c'est tout.
% archive sheet number: 17
